\documentclass[10pt,a4paper]{amsart}
\usepackage[margin=19mm]{geometry}
\usepackage{amsmath,amssymb,mathtools}
\usepackage[scr=rsfs,cal=euler]{mathalfa}
\usepackage{xfrac}
\usepackage[unicode,hidelinks,bookmarks=false]{hyperref}
\newcommand{\ie}{i.e.\ }
\newcommand{\cf}{cf.\ }
\newcommand{\resp}{respectively\ }
\newcommand{\Subsection}{\subsection}
\providecommand{\nobreakdash}{\nobreak\hbox{-}}
\setlength{\emergencystretch}{2em}
\multlinegap0pt
\usepackage{mathtools}
\usepackage[all]{xy}
\usepackage{amssymb}
\usepackage{tikz}
\newtheorem{lemma}{Lemma}
\newcommand{\R}{{\mathbb R}}
\newcommand{\rConf}{\operatorname{rConf}}
\title{Intrinsic convergence of the homological Taylor tower for $r$-immersions in $\mathbb R^n$}
\author{Gregory Arone, Franjo Sarcevic}
\date{Source: arXiv:2301.03131v2 (CC BY 4.0)}
\begin{document}
\maketitle

\noindent\emph{Source and licence.} Intrinsic convergence of the homological Taylor tower for $r$-immersions in $\mathbb R^n$. Version arXiv:2301.03131v2 (CC BY 4.0), distributed by arXiv under the Creative Commons Attribution 4.0 International licence, http://creativecommons.org/licenses/by/4.0/, as recorded in the arXiv licence metadata for this exact version and shown on the abstract page https://arxiv.org/abs/2301.03131v2. Full licence notice: \texttt{licenses/arxiv-2301.03131-CC-BY-4.0.txt}.\par\smallskip

\noindent\textit{Editorial scope.} Counted unit: the article's lemma retractive (source0.tex lines 1234--1236). The article's complete original proof of it is delivered with it (source0.tex lines 1237--1278, one \texttt{proof} environment of 42 lines). No mathematics is written, repaired or re-derived: the statement and the proof are the article's own lines, in the article's own notation, and no retained line is substituted or re-flowed.  No further source-local context is needed: every cross-reference of every retained line resolves inside the two delivered spans.  Nothing was omitted from any retained span, so the package carries no omission record.  No retained line contains a citation, so the wrapper carries no bibliography and no bibliography entry.  No retained line contains a figure environment, an image-inclusion command or drawing code, so no image is delivered and the package declares no asset dependency.  Editorial formatting only, in addition: an AMS \texttt{amsart} preamble that restates the article's own declarations and macros used by the retained lines, namely the environments \texttt{lemma} on separate counters, together with the standard packages those lines need.  The retained environments print in this document as Lemma 1 in order of appearance, whereas the article prints the same items with the numbering of its own full document.  The complete native source of the article as distributed in the arXiv e-print for this exact version is bundled unchanged as \texttt{sources/135853/source0.tex}.
\begin{lemma}\label{lemma: config retractive}
The $k$-cube of spaces $$S \mapsto \rConf(\underline{k} \setminus S, \R^n) $$ is retractive for $n\ge 2$.
\end{lemma}

\begin{proof}
Let $T$ be a finite set and suppose $x\notin T$. Our first step it to construct a section to the restriction map
\[
r_{T\cup\{x\}, T}\colon \rConf(T\cup\{x\}, \R^n)\to \rConf(T, \R^n).
\]
Let $p_1\colon \R^n \to \R$ be projection onto the first coordinate. Define a map
\[
s_{T, T\cup\{x\}}\colon \rConf(T, \R^n) \to \rConf(T\cup\{x\}, \R^n)
\]
as follows. An element of $\rConf(T, \R^n)$ is a function $f\colon T\to \R^n$ with the property that no $r$ points of $T$ go to the same point. Extend $f$ to a function from $T\cup\{x\}$ by sending $x$ to
\[
(\max \{p_1f(t)\mid t\in T\}+1, 0, \ldots, 0).
\]
In words, $x$ is sent to the point of $\R^n$ whose first  coordinate is one more than the maximal first coordinate of the existing points, and all other coordinates are zero. It is clear that the image of $x$ is different from all the other points in the configuration. Thus if $f$ was an $r$-immersion, then the resulting map $T\cup\{x\}\to \R^n$ is still an $r$-immersion. We have defined a map $s_{T, T\cup\{x\}}\colon\rConf(T, \R^n) \to \rConf(T\cup\{x\}, \R^n)$.
It is clear that the following composition is the identity (not even just homotopic to the identity but is the actual identity map)
\[
\rConf(T, \R^n) \xrightarrow{s_{T, T\cup\{x\}}} \rConf(T\cup\{x\}, \R^n)\xrightarrow{r_{T\cup\{x\}, T}}\rConf(T, \R^n).
\]
It follows that $s_{T, T\cup\{x\}}$ is a section of $r_{T\cup\{x\}, T}$. Next, we need to show that whenever $x, y\notin T$, the following diagram commutes up to homotopy
$$
\begin{aligned}
\xymatrix{
\rConf(T,\R^n) \ar[r] \ar[d] & \rConf(T\cup \{x\},\R^n) \ar[d] \\
\rConf(T\cup\{y\},\R^n) \ar[r]       &  \rConf(T\cup\{x,y\},\R^n)
}
\end{aligned}
$$
It is for this step that we need to assume $n\ge 2$. Let $f\colon T\to \R^n$ represent an element of $\rConf(T,\R^n)$. The images of $f$ in $\rConf(T\cup\{x,y\},\R^n)$ under the two ways around the diagram are two extensions of $f$ from $T$ to $T\cup\{x, y\}$. One of the extensions sends $x$ to $(\max \{p_1f(t)\mid t\in T\}+1, 0, \ldots, 0)$, and sends $y$ to $(\max \{p_1f(t)\mid t\in T\}+2, 0, \ldots, 0)$. The other extension does the same thing, with $x$ and $y$ switched. It is clear that one can write a homotopy between the two maps, by swapping the images of $x$ and $y$ along a circle in the plane spanned by the first two coordinates of $\R^n$.

Finally we need to check that the following mixed square commutes up to homotopy
$$
\begin{aligned}
\xymatrix{
\rConf(T \cup\{x\},\R^n) \ar[r] \ar[d] & \rConf(T,\R^n) \ar[d] \\
\rConf(T\cup\{x, y\},\R^n) \ar[r]       &  \rConf(T\cup\{y\},\R^n)
}
\end{aligned}.
$$
This, too, is clear. In fact, it is easy to check that there is a well-defined straight line homotopy between the two maps around the square.

We have shown that the section maps that we have defined make the cube of $r$-configuration spaces and restriction maps between them into a retractive cube.
\end{proof}

\end{document}
